\documentclass[a4paper,12pt]{article}
\usepackage[numbers]{natbib}
\usepackage{graphicx}
\usepackage{amsmath, amssymb}
\usepackage[utf8]{inputenc}
\usepackage[T1]{fontenc}
\usepackage{mathptmx}
\usepackage{setspace}
\usepackage{booktabs}
\usepackage{array}
\usepackage{xcolor}
\usepackage{caption}
\usepackage{subcaption}
\usepackage{float}
\usepackage{enumitem}
\usepackage{titlesec}
\usepackage[most]{tcolorbox}
\usepackage{algorithm}
\usepackage{algpseudocode}
\usepackage{mdframed}
\usepackage{colortbl}
\usepackage{placeins}
\usepackage{hyperref}
\usepackage{indentfirst}
\usepackage{longtable}
\usepackage{microtype}
\usepackage{geometry}
\usepackage{tikz}
\usetikzlibrary{arrows.meta, positioning, shapes.geometric, calc}
\onehalfspacing
\setlength{\emergencystretch}{3em}

\definecolor{PrimaryBlue}{HTML}{17324D}
\definecolor{SlateText}{HTML}{334155}
\definecolor{LightBlue}{HTML}{EEF5FB}
\definecolor{SoftGray}{HTML}{F6F8FA}

\hypersetup{
    colorlinks=true,
    linkcolor=PrimaryBlue,
    citecolor=PrimaryBlue,
    urlcolor=PrimaryBlue
}

\geometry{
    a4paper,
    left=2.5cm,
    right=2.5cm,
    top=2.5cm,
    bottom=2.5cm
}
\titleformat{\section}
    {\Large\bfseries\color{PrimaryBlue}}
    {\thesection}{0.75em}{}
    [\vspace{0.2em}\titlerule]
\titleformat{\subsection}
    {\large\bfseries\color{SlateText}}
    {\thesubsection}{0.75em}{}
\titleformat{\subsubsection}
    {\normalsize\bfseries\color{SlateText}}
    {\thesubsubsection}{0.75em}{}
\titlespacing*{\section}{0pt}{1.2em}{0.8em}
\titlespacing*{\subsection}{0pt}{0.9em}{0.45em}
\titlespacing*{\subsubsection}{0pt}{0.75em}{0.35em}

\captionsetup{
    font=small,
    labelfont={bf,color=PrimaryBlue},
    textfont=it,
    justification=centering,
    singlelinecheck=false
}

\setlist[itemize]{leftmargin=1.5em, itemsep=0.25em, topsep=0.25em}
\setlist[enumerate]{leftmargin=1.8em, itemsep=0.25em, topsep=0.25em}

\newcolumntype{L}[1]{>{\raggedright\arraybackslash}p{#1}}
\newcommand{\placeholder}[1]{\textcolor{PrimaryBlue}{[#1]}}
\newcommand{\dashboardfigure}[2]{%
    \IfFileExists{#1}{%
        \includegraphics[width=0.96\textwidth]{#1}%
    }{%
        \fcolorbox{PrimaryBlue}{SoftGray}{%
            \parbox[c][0.34\textheight][c]{0.9\textwidth}{%
                \centering\color{SlateText}\Large Dashboard image placeholder\\[0.5em]
                \normalsize Place file at \texttt{\detokenize{#1}}%
            }%
        }%
    }%
}

\newtcolorbox{insightbox}[1]{
    enhanced,
    breakable,
    colback=LightBlue,
    colframe=PrimaryBlue,
    coltitle=white,
    title=#1,
    fonttitle=\bfseries,
    boxrule=0.7pt,
    arc=2mm,
    left=2mm,
    right=2mm,
    top=1mm,
    bottom=1mm
}

\begin{document}
\pagenumbering{gobble}
\begin{titlepage}

\begin{tikzpicture}[remember picture,overlay]
    % Nét ngoài đậm (3pt) - Cách đều các mép 2cm
    \draw[line width=3pt] 
        ($(current page.north west) + (2cm,-2cm)$) 
        rectangle 
        ($(current page.south east) + (-2cm,2cm)$);
        
    % Nét trong thanh (0.5pt) - Cách đều các mép 2.1cm
    \draw[line width=0.5pt] 
        ($(current page.north west) + (2.1cm,-2.1cm)$) 
        rectangle 
        ($(current page.south east) + (-2.1cm,2.1cm)$);
\end{tikzpicture}

\begin{center}

%====================================
% TÊN TRƯỜNG & VIỆN
%====================================
\vspace*{-0.5cm} % Dùng số âm (-0.5cm) để kéo khối chữ ngược lên sát viền trên hơn
{\Large \textbf{HANOI UNIVERSITY OF SCIENCE AND TECHNOLOGY}}\\[0.2cm]
{\large \textbf{School of Information and Communications Technology}}

\vspace{0.8cm} % Đã giảm từ 1.2cm xuống 0.8cm để tiết kiệm không gian

\vspace{0.8cm} % Đã giảm từ 1.2cm xuống 0.8cm để tiết kiệm không gian

%====================================
% LOGO
%====================================
\includegraphics[width=0.30\textwidth]{Soict_hust_logo.jpeg}
\vspace{0.8cm} % Đã giảm từ 1.2cm xuống 0.8cm

%====================================
% LOẠI BÁO CÁO & MÔN HỌC
%====================================
{\LARGE \textbf{PROJECT REPORT}}\\[0.4cm]
{\Large \textbf{Course: Data Visualization}}

\vspace{1cm} 

{\Large \textit{Topic:}}\\[0.3cm]
{\LARGE \textbf{Global CO\textsubscript{2} Emissions:}}\\[0.2cm]
{\Large \textbf{Data Analysis with Power BI}}

\vspace{1.5cm} % Ép một khoảng trống cố định tối thiểu để đẩy tách phần Supervisor ra
\vfill 

%====================================
% THÔNG TIN GVHD & SVTH (GOM KHỐI CANH TRÁI)
%====================================
\end{center} % Kết thúc căn giữa để chuyển sang khối căn lề trái

\begin{flushleft}
\hspace*{2.5cm} % Đẩy toàn bộ khối lùi vào 2.5cm cho cân đối với khung
\begin{minipage}{12cm} % Dùng minipage để khóa chặt bố cục, tránh xô lệch
\normalsize
\renewcommand{\arraystretch}{1.3}

\textbf{Supervisor:}\quad \textbf{Dr. Tran Viet Trung} \\[0.4cm]
\textbf{Student Members:} 
\vspace{0.2cm} % Khoảng cách từ chữ Student Members xuống bảng được căn chỉnh lại

% Bảng danh sách sinh viên
\hspace*{0.5cm} 
\begin{tabular}{r @{\hspace{0.3cm}} l @{\hspace{0.8cm}} l @{\hspace{0.8cm}} l}
1. & Tran Dang Minh     & 202416720 & Data Science \& AI \\
2. & Nguyen Hoang Quan  & 202416738 & Data Science \& AI \\
3. & Le Minh Duc        & 202416677 & Data Science \& AI \\
4. & Vu Ha Anh Duc      & 202416675 & Data Science \& AI \\
5. & Pham Minh Hieu     & 202416693 & Data Science \& AI \\
\end{tabular}

\vspace{0.3cm} % Khoảng cách từ bảng xuống chữ Class ID (bù trừ để đối xứng với phần trên)

\textbf{Class ID:}\quad 168068 \\[0.1cm]
\textbf{Semester:}\quad 20252 

\end{minipage}
\end{flushleft}

\vspace*{\fill} % Tự động lấp đầy khoảng trống phía trên

\begin{center}
%====================================
% ĐỊA ĐIỂM & THỜI GIAN
%====================================
{\large \textbf{Hanoi, June 2026}}
\end{center}

\vspace*{\fill} % Tự động lấp đầy khoảng trống phía dưới với tỷ lệ 1:1 so với bên trên

\end{titlepage}

\clearpage
\pagenumbering{arabic}
\renewcommand{\contentsname}{Table of Contents}
\tableofcontents

\newpage
\section{Introduction}
Since the onset of industrialisation in the mid-eighteenth century, human activity has steadily transferred billions of tonnes of carbon from the Earth's crust into the atmosphere. 
As a result, carbon dioxide emissions have become the most pressing environmental challenge of our time, fundamentally altering the global climate system. Today, the world emits
approximately 37.8 gigatonnes of CO\textsubscript{2} annually, contributing directly to an estimated
1.3°C of global warming above pre-industrial levels. However, these aggregate statistics obscure a significant structural truth: emissions are neither uniformly distributed among countries nor constant over time. China alone constitutes around one-third of the present yearly output, whereas other Gulf states exhibit per-capita footprints that far surpass those of the world's most impoverished places. Understanding how these disparities arose, and how responsibility
should be apportioned between early industrialisers and rapidly developing economies, is
inseparable from any credible discussion of climate policy.
Despite the complexity of this global picture, emissions data are frequently encountered as figures embedded in news reports, policy briefs, or static tables, leaving
analysts without a unified, longitudinal view that integrates national totals, per-capita
measures, sectoral breakdowns, and historical trajectories in a single analytical framework.
To address this gap, this project draws on the Our World in Data CO\textsubscript{2} and Greenhouse Gas
Emissions dataset, openly maintained at \texttt{github.com/owid/co2-data}. This collection offers systematically organized, annually revised data encompassing over two centuries of emissions across nations, regions, and income categories, with detailed breakdowns by fuel type, land use, and consumption-based accounting. Although the dataset has a comprehensive quantitative record, these records exist as individual observations lacking systematic visualization. The main goal of this project is to convert this extensive longitudinal dataset into organized, actionable climate insights via an interactive analytical dashboard.

To achieve this, we developed a multi-page Streamlit dashboard designed to consolidate
emissions records into one coherent analytical environment, revealing patterns by time
period, geography, income group, and fuel source. The dashboard is structurally divided
into three main analytical pages:

\begin{itemize}
    \item \textbf{Page 1 --- Global Snapshot:} provides a macro-level summary addressing
    critical questions about the current scale of emissions, which nations bear the greatest
    absolute and per-capita burden, and how geographic concentration has evolved across
    hemispheres and income groups.

    \item \textbf{Page 2 --- Historical Evolution:} shifts to a longitudinal perspective,
    tracing global CO\textsubscript{2} trajectories from 1750 to the present. Annotated timeline charts
    expose the precise historical moments---the post-WWII industrial boom, the 1970s oil
    crisis, and the COVID-19 disruption---that shaped the emissions curve, while regional
    decompositions reveal the transfer of cumulative responsibility from early Western
    industrialisers to rapidly developing economies in Asia.

    \item \textbf{Page 3 --- Current Trend Analysis:} examines recent trajectories and
    decarbonisation progress, contrasting nations on accelerating emissions pathways against
    those achieving measurable reductions, and disaggregating China's evolving fuel mix to
    illustrate the complexities of large-economy energy transition.
\end{itemize}

A core analytical requirement of this project is distinguishing between two complementary
dimensions of climate impact: absolute emissions volume---measured in gigatonnes of
CO\textsubscript{2} and relevant for assessing total atmospheric loading---and relative per-capita
intensity, which exposes structural inequalities in how emissions are distributed across
populations and income levels. Rather than merely reporting aggregate figures, the dashboard
supports comparative and exploratory analysis for target audiences including climate policy
researchers, sustainability analysts, educators, and data journalists.

Furthermore, this report documents the data visualisation rationale, visual encoding
strategies, and analytical narrative underpinning each dashboard component, while addressing
key limitations inherent to the dataset---including reporting lags in developing-economy
statistics, the non-trivial distinction between production-based and consumption-based
accounting, the aggregation of politically complex entities within regional groupings, and
the interpretive caution required when attributing causal significance to observed
correlations between economic development indicators and emissions trajectories.

\section{Dataset Description}
The project makes use of the Our World in Data CO\textsubscript{2} and Greenhouse Gas Emissions dataset, which monitors regional and national carbon dioxide emissions in over 200+ countries and territories from 1750 to the present. Every record is an annual observation for a particular nation or collective entity (such as a continent or income group). Table~\ref{tab:dataset-variables} provides a summary of the main variables used in the Streamlit dashboard.

\begin{table}[H]
\centering
\caption{Main variables used from the OWID CO\textsubscript{2} and Greenhouse Gas Emissions Dataset.}
\label{tab:dataset-variables}
\renewcommand{\arraystretch}{1.25}
\begin{tabular}{L{0.32\textwidth} L{0.57\textwidth}}
\toprule
\textbf{Variable} & \textbf{Description} \\
\midrule
\texttt{country}            & Name of the country, region, or aggregate entity
                              (e.g., continent or income group). \\
\texttt{year}               & Year of the observation, ranging from 1750 to 2024. \\
\texttt{iso\_code}          & ISO 3166-1 alpha-3 country code used for geographic
                              mapping and cross-referencing. \\
\texttt{co2}                & Total annual CO\textsubscript{2} emissions from fossil fuels and
                              industry, measured in million tonnes. \\
\texttt{co2\_per\_capita}   & Annual CO\textsubscript{2} emissions per person, measured in
                              tonnes per capita. \\
\texttt{cumulative\_co2}    & Cumulative CO\textsubscript{2} emissions since 1750, measured
                              in million tonnes. \\
\texttt{co2\_per\_gdp}      & CO\textsubscript{2} emissions per unit of GDP (carbon
                              intensity), measured in kg per international USD. \\
\texttt{coal\_co2}          & CO\textsubscript{2} emissions attributable to coal
                              combustion, measured in million tonnes. \\
\texttt{oil\_co2}           & CO\textsubscript{2} emissions attributable to oil
                              combustion, measured in million tonnes. \\
\texttt{gas\_co2}           & CO\textsubscript{2} emissions attributable to natural gas
                              combustion, measured in million tonnes. \\
\texttt{population}         & Total population of the country or entity in the
                              given year. \\
\texttt{gdp}                & Gross domestic product in international USD, used as
                              an economic development indicator. \\
\bottomrule
\end{tabular}
\normalsize
\setlength{\tabcolsep}{6pt}
\end{table}

There are two factors that are particularly crucial for interpretation. The variable \texttt{co2} represents a nation's direct contribution to overall atmospheric loading by measuring the absolute scale of its yearly emissions. \texttt{co2\_per\_capita}, on the other hand, calculates the relative intensity of emissions in relation to the size of the country's population.
Because of the size of its population, a huge, highly industrialized economy may release hundreds of millions of tons annually despite registering a relatively low per-capita statistic.
On the other hand, a small, wealthy petrostate may emit significantly less tonnes overall, but its per-capita footprint may be among the highest in the world, revealing a structural reliance on carbon-intensive activities.
The dashboard may compare nations based on their historical responsibility—that is, their entire contribution to atmospheric CO2 concentration since industrialization started—rather than just their present yearly output thanks to the variable \texttt{cumulative\_co2}. This lends credence to the examination of whether early industrializers in North America and Europe had a disproportionate amount of cumulative responsibility in relation to their current emissions. Additionally, the fuel-disaggregated variables \texttt{coal\_co2}, \texttt{oil\_co2}, and \texttt{gas\_co2} offer crucial background information about each nation's energy mix, enabling the dashboard to track structural changes in the composition of emissions over time and pinpoint economies that continue to rely significantly on coal combustion as a major source of decarbonization risk.

\section{Data Preprocessing}

The preprocessing stage transformed the original Our World in Data (OWID) CO\textsubscript{2}
and Greenhouse Gas Emissions file into a set of dashboard-ready analytical tables. This
stage was implemented primarily in the modules under \texttt{src/data/}, supported by the
exploratory audit and cleaning notebooks \texttt{01\_data\_audit.ipynb},
\texttt{02\_cleaning\_feature\_engineering.ipynb}, and
\texttt{03\_eda\_insight\_discovery.ipynb}. The purpose of preprocessing was not only to
remove technical inconsistencies, but also to restructure the dataset according to the
analytical needs of the dashboard: country comparison, world and regional historical
trends, cumulative responsibility, recent emissions momentum, and fuel-source
decomposition.

\subsection{Raw Dataset Overview}

The raw input file, \texttt{data/raw/owid-co2-data.csv}, is the OWID CO\textsubscript{2}
dataset in country--year format. Each row represents one annual observation for a named
entity and year. In the version used for this project, the raw file contains 50,411 rows
and 79 columns, covering the period from 1750 to 2024. The \texttt{country} field includes
both sovereign countries and non-country aggregate entities such as \texttt{World},
continents, income groups, international aviation, and international shipping. The
dataset contains 254 distinct named entities before preprocessing.

Although the OWID file provides a broad climate and energy accounting record, the dashboard
does not use all columns. The project focuses on variables directly relevant to the three
dashboard pages: total CO\textsubscript{2} emissions (\texttt{co2}), emissions per capita
(\texttt{co2\_per\_capita}), cumulative emissions (\texttt{cumulative\_co2}), population,
GDP, and fuel-specific emissions such as \texttt{coal\_co2}, \texttt{oil\_co2},
\texttt{gas\_co2}, \texttt{cement\_co2}, \texttt{flaring\_co2}, and
\texttt{land\_use\_change\_co2}. These variables were selected because they map directly
to the intended visual questions: current scale, per-capita inequality, historical
responsibility, regional differences, recent trend direction, and fuel-mix composition.
Many remaining OWID columns were not used because they either represented alternative
accounting systems outside the scope of the dashboard or had very high structural
missingness. For example, in the raw audit, variables such as
\texttt{share\_global\_cumulative\_other\_co2}, \texttt{share\_global\_other\_co2},
\texttt{other\_co2\_per\_capita}, and several consumption-based variables were missing
for more than 89\% of records. Excluding these fields avoided cluttering the dashboard
with sparse indicators that would not support stable cross-country comparison.

\subsection{Data Audit}

The initial audit notebook examined the raw dataset's dimensions, data types, missing
values, duplicate records, year coverage, and country/entity coverage. The raw file was
confirmed to contain no duplicate \texttt{country}--\texttt{year} pairs. The year range
was validated as 1750--2024, which supports both long-run historical analysis and
latest-year snapshot views.

Missingness was assessed by column rather than handled through broad imputation. This was
important because the OWID dataset is historically heterogeneous: early years contain
fewer variables, some fuel-specific series begin later than total CO\textsubscript{2}, and
GDP indicators are not available for every country-year. The audit therefore treated
missingness as an availability characteristic of the source data rather than as a simple
data-entry error. For example, in the latest year table generated after preprocessing,
\texttt{co2} is available for 216 out of 219 country rows, \texttt{co2\_per\_capita} for
214 out of 219 rows, and \texttt{population} for 217 out of 219 rows, while \texttt{gdp}
is not available in the 2024 records. This limitation is handled downstream by using the
latest available GDP within the relevant recent window when GDP per capita is required
for trend analysis.

A central part of the audit was identifying the difference between true country rows and
aggregate rows. This distinction is essential because the raw dataset stores countries,
continents, world totals, and income-group totals in the same table. If these rows were
ranked together, a chart of top emitters would incorrectly compare individual countries
against entities such as \texttt{World} or \texttt{Asia}. Conversely, removing all
aggregates would make it impossible to produce accurate world and regional historical
trend lines. The preprocessing pipeline therefore explicitly separates country-level data
from aggregate entities while preserving both for their appropriate analytical purposes.

\subsection{Cleaning Process}

The cleaning process is implemented in \texttt{src/data/clean\_data.py}. The first step
standardizes all column names into lowercase snake case using a regular-expression based
normalization function. This ensures that downstream code can refer consistently to fields
such as \texttt{co2\_per\_capita}, \texttt{cumulative\_co2}, and \texttt{land\_use\_change\_co2}
without depending on the raw file's original naming format.

The pipeline then validates the presence of the two required structural fields,
\texttt{country} and \texttt{year}. Rows with missing values in either field are removed,
because they cannot be placed in a country--year time series. The \texttt{year} column is
converted to integer type after coercion, while metric columns listed in
\texttt{METRIC\_COLUMNS} are converted to numeric values. Non-negative metrics such as
\texttt{co2}, \texttt{co2\_per\_capita}, \texttt{population}, \texttt{gdp}, and the main
fuel-specific CO\textsubscript{2} variables are checked for negative values; invalid
negative entries are converted to missing values. The exception is
\texttt{land\_use\_change\_co2}, which is allowed to be negative because land-use change
can represent a net sink in some years.

The pipeline does not fill missing emissions values with zero by default. Missing values
are preserved as missing, and individual charts or analytical functions later drop records
only when the metric required for that visual is absent. This design reduces the risk of
misrepresenting unreported values as true zeros.

Country and aggregate separation is performed through an \texttt{is\_aggregate} flag. A
row is treated as an aggregate if its \texttt{country} value belongs to a known set of
non-country entities, if the entity is one of the six dashboard regions, or if its
\texttt{iso\_code} starts with the OWID aggregate prefix \texttt{OWID\_}. Rows with
invalid or missing ISO codes are also treated cautiously unless they appear in the
country-region override table. As a result, country rankings, top-emitter bar charts, and
country-level scatter plots are computed only from country rows, while world and regional
trend views use the preserved aggregate rows.

\subsection{Feature Engineering}

Feature engineering is implemented in \texttt{src/data/feature\_engineering.py} and
combined in \texttt{src/data/build\_datasets.py}. First, the cleaning step creates
\texttt{region} and \texttt{income\_group} fields. Region assignment is inferred from the
raw \texttt{continent} field when available, from explicit region names for regional
aggregate rows, and from a manually maintained country-to-region override table. The code
also includes an optional ISO-based enrichment step using \texttt{country\_converter},
which fills missing region labels from ISO3 metadata when possible. After the most recent
rebuild, the latest-year country table contains no \texttt{Unknown} region labels. The
\texttt{income\_group} field is more limited: the current raw OWID file does not provide
complete country-level income classifications, so the pipeline mainly preserves income
aggregate entities such as high-income and low-income country groups. This is a limitation
rather than an imputed classification.

The pipeline computes \texttt{share\_of\_world\_co2} by joining every row to the
\texttt{World} row for the same year and calculating:
\[
    \texttt{share\_of\_world\_co2}
    =
    \frac{\texttt{co2}}{\texttt{World co2 in the same year}} \times 100.
\]
This field supports map tooltips, KPI context, and country ranking interpretation. A
similar method is used to create \texttt{cumulative\_share\_of\_world} from
\texttt{cumulative\_co2}, allowing the dashboard to compare historical responsibility as
a share of global cumulative emissions.

The project also creates temporal trend features. Previous-year values are added for
\texttt{co2}, \texttt{co2\_per\_capita}, and \texttt{population}. For five-year momentum,
the code computes \texttt{co2\_start\_5y}, \texttt{co2\_change\_5y}, and
\texttt{co2\_cagr\_5y}. The five-year compound annual growth rate is calculated only when
both the current and start values are positive:
\[
    \texttt{co2\_cagr\_5y}
    =
    \left(\frac{\texttt{co2}_t}{\texttt{co2}_{t-5}}\right)^{1/5} - 1.
\]
For the current trend scatter plot, the analysis module can also compute a custom CAGR
between selected years. In the present dashboard configuration, the recent quadrant view
uses 2022--2024 CO\textsubscript{2} CAGR and GDP per capita derived from the latest
available GDP and population values within that window. This is necessary because GDP is
not available in the 2024 country rows of the current OWID file.

Fuel-source decomposition is prepared by reshaping the country-year table from wide to
long format. The columns \texttt{coal\_co2}, \texttt{oil\_co2}, \texttt{gas\_co2},
\texttt{cement\_co2}, \texttt{flaring\_co2}, and \texttt{land\_use\_change\_co2} are
melted into a pair of fields: \texttt{fuel\_source} and \texttt{value}. This long format
is required for stacked area charts and regional fuel-mix decomposition, because each fuel
source becomes a categorical series rather than a separate column.

\subsection{Output Datasets}

The build script writes a set of processed and aggregated CSV files to
\texttt{data/processed/} and \texttt{data/aggregated/}. Table~\ref{tab:preprocessing-outputs}
summarizes the generated files after the latest rebuild.

\begin{longtable}{L{0.31\textwidth} L{0.14\textwidth} L{0.45\textwidth}}
\caption{Dashboard-ready datasets generated by the preprocessing pipeline.}
\label{tab:preprocessing-outputs}\\
\toprule
\textbf{Output file} & \textbf{Shape} & \textbf{Purpose in the dashboard} \\
\midrule
\endfirsthead
\toprule
\textbf{Output file} & \textbf{Shape} & \textbf{Purpose in the dashboard} \\
\midrule
\endhead
\texttt{co2\_cleaned.csv} & 50,411 $\times$ 90 &
Complete cleaned dataset containing both countries and aggregate entities, with standardized
columns and engineered features. \\
\texttt{co2\_country\_year.csv} & 42,654 $\times$ 90 &
Country-level panel used for country rankings, maps, treemaps, scatter plots, and
country-specific KPI calculations. Aggregate rows are excluded from this table. \\
\texttt{co2\_aggregates.csv} & 7,757 $\times$ 90 &
Aggregate entities such as \texttt{World}, continents, and income groups. Used where
regional or global totals are analytically appropriate. \\
\texttt{co2\_latest\_year.csv} & 219 $\times$ 90 &
Latest-year country snapshot, used for current-state comparisons and KPI-style summary
views. \\
\texttt{co2\_trends.csv} & 1,314 $\times$ 90 &
Recent country-year subset covering the latest five-year window in the rebuilt dataset,
supporting momentum analysis and trend-oriented views. \\
\texttt{co2\_fuel\_long.csv} & 255,924 $\times$ 7 &
Long-format fuel-source emissions table used for stacked area charts and fuel-mix
decomposition by region. \\
\texttt{global\_yearly.csv} & 275 $\times$ 90 &
World-level annual time series used for the historical global emissions line chart. \\
\texttt{regional\_yearly.csv} & 1,650 $\times$ 90 &
Annual time series for Asia, Europe, North America, South America, Africa, and Oceania,
used for regional comparison charts. \\
\texttt{country\_latest\_metrics.csv} & 219 $\times$ 90 &
Aggregated latest-year country metrics for snapshot cards, ranking, and map-oriented
views. \\
\bottomrule
\end{longtable}

\subsection{Relevance to Visualization Design}

The preprocessing pipeline directly supports the visual storytelling goals of the
dashboard. The Global Snapshot page requires a clean latest-year country table for KPI
cards, choropleth mapping, treemap composition, and top-country ranking. The Historical
Evolution page requires world and regional aggregate rows to be preserved, because the
global trajectory and continental comparisons cannot be reconstructed reliably from
country rows alone without additional aggregation assumptions. The Current Trend Analysis
page requires recent country-level trend features, GDP-per-capita preparation, and a
long-format fuel table to support quadrant analysis and fuel-source decomposition.

Separating country rows from aggregate rows is especially critical for visual accuracy.
If aggregate rows were left inside country rankings, \texttt{World}, \texttt{Asia}, or
income-group totals would dominate charts intended to compare individual countries. This
would distort the analytical message and produce misleading conclusions. Conversely,
discarding aggregates would weaken the historical narrative by removing authoritative
world and regional totals. The project therefore uses two complementary data streams:
country-level tables for comparison and ranking, and aggregate tables for global and
regional historical context.

The engineered features also improve interaction and cross-filtering. Region labels make
it possible to filter and color charts consistently by continent. World-share indicators
provide interpretable tooltip context for map and ranking views. Cumulative responsibility
features support the distinction between current annual emissions and historical
contribution. CAGR and absolute-change fields enable the dashboard to identify countries
whose emissions are accelerating or declining. Finally, the long fuel-source table allows
fuel composition to be visualized as a time-varying stacked structure rather than as a
set of disconnected columns. In this sense, preprocessing is not a preliminary technical
step only; it defines the analytical grammar that makes the dashboard coherent,
interactive, and suitable for climate-emissions storytelling.


\section{Dashboard Design Overview}
With the dataset fully cleaned and structural missingness accounted for, the next step is to visually encode these records into an interactive format. The Power BI dashboard is organized into two analytical pages. Page 1 is designed for a macro-level overview of global trends, while Page 2 focuses on a micro-level perspective regarding company maturity and funding.

\subsection{Page 1: Global Layoff Overview}
Page 1 provides a high-level summary of global workforce reductions. It is designed to answer the core question: \textit{When, in which industries, and in which geographic regions did layoffs occur most severely?}

The page incorporates the following visual components:

\begin{itemize}
    \item \textbf{KPI Cards (Total Employees Laid off, Peak Monthly Layoff, Total Industries)}: These cards provide an immediate statistical snapshot before the viewer examines detailed charts. The \textit{Peak Monthly Layoff} card specifically directs attention to the most critical period in the temporal trend.
    \item \textbf{Date Slicer}: Allows the viewer to dynamically filter layoff records by specific time periods.
    \item \textbf{Area Chart (Total Layoffs over Time)}: Visualizes the temporal trend of layoffs, effectively highlighting periods of sudden spikes or declines over time.
    \item \textbf{Horizontal Bar Chart (Total Layoffs by Industry)}: Compares the absolute scale of layoffs across different sectors, quickly identifying the hardest-hit industries.
    \item \textbf{Choropleth Map (Total Layoffs by Country)}: Maps geographic distribution using color intensity (saturation). It utilizes tooltips to display additional context, such as the total number of companies affected in a specific region.
\end{itemize}

Together, these views support a compact yet multidimensional overview. The KPI cards establish the baseline scale, the area chart emphasizes temporal movement, the bar chart categorizes industry vulnerability, and the map expands the analysis into geographic space.

\begin{figure}[H]
\centering
\dashboardfigure{imgs/LayoffDV-1.pdf}{Page 1 dashboard placeholder}
\caption{Overview dashboard summarizing global layoffs across time, industry, and geography. KPI cards provide aggregate measures, while the area chart, bar chart, and shape map support temporal, categorical, and spatial analysis.}
\label{fig:page1-overview}
\end{figure}

\subsection{Page 2: Company \& Funding Stage Analysis}
Page 2 shifts the focus to structural firm characteristics, examining the relationship between funding stage and layoff severity. It is designed to answer: \textit{How does a company's funding stage relate to both the percentage and absolute scale of its layoffs?}

The page contains the following visual components:
\begin{itemize}
    \item \textbf{KPI Cards (Total Employees Laid off, Total Funds Raised, Total Companies)}: Summarize the scale, accumulated capital, and scope of companies under the selected filters.
    \item \textbf{Date Slicer}: Enables temporal filtering specific to the funding-stage analysis.
    \item \textbf{Vertical Bar Chart (Mean Layoff Percentage by Funding Stage)}: Compares the relative layoff severity across different maturity stages, highlighting which stages faced the most intense internal restructuring.
    \item \textbf{Horizontal Bar Chart (Total Layoffs by Company)}: Identifies individual companies contributing the largest absolute layoff counts, preventing the analysis from being skewed solely by aggregate stage averages.
    \item \textbf{Scatter Plot (Total Funding vs. Layoff by Stage)}: Examines the correlation between accumulated funding and absolute layoff size. This reveals whether highly funded companies experienced proportionally smaller or larger layoffs.
\end{itemize}

The analytical strength of this page lies in its juxtaposition of relative and absolute metrics. One funding stage may exhibit a high \texttt{percentage\_laid\_off} (indicating severe internal disruption), while another may show a high \texttt{total\_laid\_off} (indicating broader labor-market impact). The inclusion of the company bar chart and scatter plot ensures that firm-level realities are not lost in high-level aggregations.

\begin{figure}[H]
\centering
\dashboardfigure{imgs/LayoffDV-2.pdf}{Page 2 dashboard placeholder}
\caption{Funding-stage analysis dashboard examining layoff severity by company maturity stage. KPI cards summarize selected records, while the column chart, company bar chart, and scatter plot connect funding stage, firm-level scale, and the funding--layoff relationship.}
\label{fig:page2-stage-analysis}
\end{figure}

\section{Technique Application by Chapter}
This section connects the theoretical principles of data visualization to the concrete
design choices made in the CO\textsubscript{2} emissions dashboard. The project does not
treat visualization as decoration; rather, it uses visual representation as an analytical
tool that converts a large country--year table into interpretable spatial, temporal, and
comparative patterns. The discussion below follows the major topics summarized in the
course material: the foundations of visualization, visual encoding, graphical perception,
figure design, data-specific visualization techniques, and storytelling with data.

\subsection{Chapter 1: Foundations and Philosophy of Data Visualization}

The first principle applied in this project is that visualization acts as an external
cognitive aid. The OWID CO\textsubscript{2} dataset contains tens of thousands of annual
records across countries, regions, fuels, and historical periods. In raw tabular form,
the user can inspect values, but cannot easily perceive long-run acceleration, regional
redistribution, or the contrast between total emissions and per-capita inequality.
The dashboard therefore transforms abstract symbols---country names, years, and numerical
emissions values---into geometric forms that extend human memory and pattern recognition.

The project belongs primarily to \textit{information visualization}. Although the subject
matter is environmental and scientific, the dashboard does not visualize physical climate
simulations or three-dimensional atmospheric processes. Instead, it visualizes abstract
statistical records: country-year emissions, cumulative totals, regional categories, GDP
per capita, and fuel-source shares. This distinction guided the design toward maps, line
charts, bar charts, treemaps, and scatter plots rather than scientific volume rendering or
3D climate models.

The dashboard also balances three core values of visualization. First, it records
information by presenting cleaned country-year and aggregate statistics in a stable
visual structure. Second, it supports analytical reasoning by enabling comparisons such
as China versus the United States, Asia versus Europe, or total CO\textsubscript{2} versus
CO\textsubscript{2} per capita. Third, it communicates an argument: global emissions are
historically cumulative, geographically concentrated, and structurally linked to both
population scale and energy composition. In the spirit of Anscombe's Quartet, the
dashboard assumes that summary statistics alone are insufficient; patterns such as
post-war acceleration, regional divergence, and outlier countries only become clear when
the data are visualized.

\subsection{Chapter 2: Visual Encoding and Visual Model}

The dashboard applies the visual model by mapping data attributes to marks and perceptual
channels. Country records, regions, and fuel sources become graphical marks, while
quantitative measures such as \texttt{co2}, \texttt{co2\_per\_capita},
\texttt{cumulative\_co2}, and \texttt{co2\_cagr} are encoded through position, length,
area, size, or color intensity.

Several encoding decisions follow the principle that different data types require
different visual channels. Quantitative comparison is handled mainly through position and
length, the most effective channels for numerical judgment. This is why top emitters,
cumulative responsibility, fastest-increasing countries, and fastest-declining countries
are shown as sorted horizontal bar charts. The temporal structure of the dataset is
encoded by position along the x-axis in the global and regional line charts. Nominal
categories such as region and fuel source are encoded with hue, using a colorblind-aware
qualitative palette inspired by Okabe and Ito. Ordered magnitude on the world map is
encoded with a sequential color scale, because map color intensity should express more
or less emissions rather than different categories.

\begin{itemize}
    \item \textbf{KPI cards:} use textual marks and compact summary blocks to encode the
    current global total, largest emitter, per-capita leader, and selected-country context.
    \item \textbf{Choropleth map:} uses geographic area marks; color intensity encodes
    either total CO\textsubscript{2} or CO\textsubscript{2} per capita, while tooltips
    expose additional variables such as region, population, and world share.
    \item \textbf{Treemap:} uses containment and area to represent region--country
    composition. The hierarchy \texttt{region} $\rightarrow$ \texttt{country} makes the
    part-to-whole structure visible.
    \item \textbf{Line charts:} use line marks and position to encode emissions change
    over time.
    \item \textbf{Scatter plot:} uses point marks where x-position encodes recent
    CO\textsubscript{2} CAGR, y-position encodes GDP per capita, color encodes region,
    and marker size encodes current total CO\textsubscript{2}.
    \item \textbf{Stacked area chart:} uses vertical proportion over time to encode fuel
    composition within a selected region.
\end{itemize}

These choices follow Mackinlay's criteria of expressiveness and effectiveness. The
visuals are expressive because each chart encodes only variables that are meaningful for
the analytical question. They are effective because high-priority quantitative comparisons
are assigned to position and length, while color hue is reserved for categorical grouping
and color intensity for ordered magnitude.

\subsection{Chapter 3: Graphical Perception}

The dashboard design accounts for how viewers perceive graphical marks. Since people
estimate length more accurately than area or angle, the project avoids pie charts and
uses horizontal bar charts for country ranking. This is especially important for top
emitters and cumulative responsibility, where small differences between high-ranking
countries must remain readable. Bars also begin at zero, preventing visual exaggeration
of differences in emissions magnitude.

Pre-attentive processing is used carefully rather than excessively. Darker countries on
the choropleth map immediately attract attention as high-emission areas. In the scatter
plot, the quadrant guide lines and contrasting region colors help the viewer quickly
separate countries with positive emissions growth from those with declining growth. The
dashboard avoids stacking too many simultaneous encodings in one figure; for example,
the map uses sequential color for magnitude, while region hue is mainly used in treemaps,
bar charts, line charts, and the scatter legend.

Gestalt principles also shape the layout. KPI cards are grouped as a compact summary
strip, encouraging the viewer to read them as one dashboard header. Related charts are
placed near each other: on the Global Snapshot page, the map and treemap appear side by
side because both describe spatial composition; on the Historical Evolution page, the
global line chart precedes cumulative and regional comparisons; on the Current Trend
Analysis page, the quadrant scatter and two momentum rankings form a single analytical
block. This proximity reduces the need for the viewer to remember information across
distant parts of the screen.

The design also recognizes change blindness. Long historical trends can be difficult to
interpret if viewers must infer turning points unaided, so the global line chart includes
annotations for major historical periods such as industrialization, the post-war boom,
the oil crisis, the financial crisis, and the COVID-19 disruption. These annotations
direct attention to meaningful changes without requiring animation.

\subsection{Chapter 4: Principles of Figure Design}

The project follows the principle of maximizing the data-ink ratio. Decorative elements
are minimized, and the dashboard avoids chart junk such as 3D effects, unnecessary
illustrations, and dense explanatory text inside the application. Recent interface
iterations removed redundant insight boxes so that screen space is used primarily for the
charts themselves. This improves the ratio between visual ink devoted to data and ink
devoted to decoration.

Graphical integrity is addressed by matching visual magnitude to numerical magnitude.
Bar charts use length to encode emissions and start from zero, satisfying the principle
of proportional ink. The treemap uses area for part-to-whole composition, and therefore
is reserved for hierarchical composition rather than precise ranking. The bubble scatter
uses marker size only as a secondary cue for total CO\textsubscript{2}; precise values are
provided in tooltips, acknowledging that humans judge area less accurately than position.

Axis design also reflects the course principles. The historical line charts use time on
the horizontal axis and emissions on the vertical axis, making temporal change perceptually
natural. The current-trend scatter plot fixes the x-axis around zero so that positive and
negative CO\textsubscript{2} growth are visually balanced. GDP per capita is displayed on
a transformed scale with readable monetary ticks, because economic development indicators
are highly skewed across countries. This prevents high-income outliers from compressing
lower-income countries into an unreadable band.

The dashboard avoids animation as a primary comparison mechanism. Although CO\textsubscript{2}
data are temporal, animated year-by-year changes would force the viewer to rely on memory.
Instead, the project uses static line charts, sorted rankings, and side-by-side comparisons,
which allow the eye to compare values directly. This follows the principle that, for
analytical comparison, ``eyes beat memory.''

\subsection{Chapter 5: Application to Specific Data Types}

The OWID CO\textsubscript{2} dataset is mainly tabular, but it contains several analytical
structures: temporal records, geographic entities, hierarchical region--country
relationships, and fuel-source categories. The project selects visualization techniques
according to these structures.

For table data, the dashboard uses KPI cards, sorted bar charts, and scatter plots to
show quantities, rankings, and relationships. For temporal data, it uses line charts and
stacked area charts. The global line chart emphasizes the long-run trajectory of world
emissions, while the regional line chart compares continental trajectories over the same
time period. For geographic data, the dashboard uses a choropleth map based on ISO country
codes. This is appropriate because the analysis concerns country-level spatial
distribution rather than point locations. For hierarchical data, the treemap represents
the relationship between regions and their constituent countries. For compositional fuel
data, the long-format fuel table supports a 100\% stacked area chart, making it possible
to compare how the share of coal, oil, gas, cement, flaring, and land-use change evolves
over time.

Graph and network visualization techniques are intentionally not used. The dataset does
not contain explicit node-link relationships such as trade networks, supply chains, or
country-to-country emission flows. Therefore, force-directed graphs, edge bundling, and
network layouts would add visual complexity without representing actual relationships in
the data. Similarly, sunburst charts were not selected because the dashboard requires
comparative readability and temporal analysis more than deep nested hierarchy.

\subsection{Chapter 6: Storytelling with Data}

The dashboard uses data storytelling to turn a large historical dataset into a coherent
analytical sequence. The narrative follows a broad-to-specific structure. The Global
Snapshot page begins with the current state of emissions: total global scale, country
distribution, regional composition, and top emitters. The Historical Evolution page then
adds temporal context, showing how emissions accumulated and how regional responsibility
changed over time. Finally, the Current Trend Analysis page shifts to momentum and
transition, asking which countries are growing or declining and how regional fuel mixes
are changing.

This structure reflects the three pillars of storytelling with data. The data pillar is
provided by the cleaned OWID country-year and aggregate tables. The visualization pillar
is implemented through carefully selected charts whose encodings match the analytical
task. The narrative pillar is provided by page sequencing, chart titles, annotations, and
interactive filters that guide the viewer from global overview to historical explanation
and then to current transition dynamics.

The dashboard also uses contrast to separate context from focus. Regional colors are
consistent across charts, while the selected country can be highlighted on map and treemap
views. The map uses a sequential scale for emissions magnitude, whereas region and fuel
legends use qualitative colors. This distinction reduces interference between categorical
and quantitative interpretations. Interaction supports exploratory storytelling: users
can select regions, switch between total and per-capita metrics, compare countries, and
change the region used for fuel decomposition. Thus, the dashboard combines an authored
storyline with reader-driven exploration.

\subsection{Summary of Techniques}
Table~\ref{tab:visual-techniques} summarizes the principal visualization techniques used
in the CO\textsubscript{2} dashboard and the role each technique plays in the analytical
story.

\begin{table}[H]
\centering
\caption{Summary of visualization techniques used in the CO\textsubscript{2} dashboard.}
\label{tab:visual-techniques}
\small
\renewcommand{\arraystretch}{1.25}
\begin{tabular}{L{0.23\textwidth} L{0.28\textwidth} L{0.37\textwidth}}
\toprule
\textbf{Technique} & \textbf{Main variables} & \textbf{Analytical role} \\
\midrule
KPI cards & \texttt{co2}, \texttt{co2\_per\_capita}, \texttt{share\_of\_world\_co2} &
Provide immediate summary of current global scale, largest emitters, and selected-country context. \\
Choropleth map & \texttt{iso\_code}, \texttt{country}, \texttt{co2}, \texttt{co2\_per\_capita} &
Shows geographic concentration of total or per-capita emissions through sequential color intensity. \\
Treemap & \texttt{region}, \texttt{country}, \texttt{co2} or \texttt{co2\_per\_capita} &
Represents region--country composition and part-to-whole structure. \\
Horizontal bar chart & \texttt{country}, \texttt{co2}, \texttt{cumulative\_co2}, \texttt{co2\_cagr} &
Supports accurate ranking of top emitters, cumulative responsibility, and fastest-changing countries. \\
Annotated line chart & \texttt{year}, \texttt{co2} &
Reveals long-run global historical trajectory and directs attention to major turning points. \\
Regional line chart & \texttt{year}, \texttt{region}, \texttt{co2} &
Compares regional emissions trajectories using consistent color categories. \\
Bubble scatter plot & \texttt{co2\_cagr}, \texttt{gdp\_per\_capita}, \texttt{co2}, \texttt{region} &
Compares recent emissions growth, economic profile, regional category, and current emissions scale. \\
Stacked area chart & \texttt{year}, \texttt{fuel\_source}, \texttt{value} &
Shows how regional fuel-source shares evolve over time. \\
Interactive filters and selection & \texttt{year}, \texttt{region}, \texttt{country}, \texttt{fuel\_source} &
Enable cross-filtering, metric switching, country selection, and region-specific decomposition. \\
\bottomrule
\end{tabular}
\normalsize
\setlength{\tabcolsep}{6pt}
\end{table}\normalsize 

\section{Economic Insight Analysis}
Building upon the visual and narrative techniques discussed in the previous section, the dashboard enables deep economic interpretation by connecting layoff scale, timing, industry, geography, funding stage, and company-level structure. The insights below synthesize the macro and micro-level realities extracted from the Power BI dashboard.

\subsection{Layoffs as a Signal of Macroeconomic Pressure}
When total layoffs increase sharply over time, it reflects pressure from the macroeconomic environment-such as higher interest rates, increased cost of capital, weaker consumer demand, and heavily revised growth expectations. Particularly, technology companies depend on expectations of future growth. When capital markets tighten, these firms often reduce operating costs more aggressively than firms in stable, traditional sectors. 

Peaks in the area chart highlight these periods of synchronized restructuring. Notably, the dashboard reveals a massive spike during \textbf{late 2022 to early 2023}, suggesting that the global tech ecosystem collectively adjusted its workforce size in response to an abrupt end to the pandemic-era zero-interest-rate environment.

\subsection{Industry-level Vulnerability}
The industry bar chart demonstrates that layoffs are not evenly distributed across sectors. Industries with the highest total layoffs are typically those that expanded aggressively during the preceding growth period and faced severe corrections when business models could no longer sustain those expansion rates.

The dashboard indicates that \textbf{Retail} (closely followed by Consumer and Hardware) contributes the largest absolute share of layoffs. This suggests that e-commerce and consumer-facing tech were heavily exposed to post-pandemic corrections and cost restructuring. This sectoral concentration proves that workforce reductions were structurally driven rather than randomly distributed.

\subsection{Geographic Concentration of Layoffs}
The choropleth map illustrates the geographic footprint of the layoffs. However, a country with high total layoffs could represent a massive tech ecosystem with thousands of firms, or merely a few dominant companies executing massive cuts.

The tooltip detailing \textit{Total Companies} contextualizes this. The \textbf{United States} emerges as the most intensely shaded region, accompanied by a dominant company count. This combination indicates systemic exposure across the entire US tech sector, rather than isolated events by a few megacorporations, cementing its position as the epicenter of the layoff wave.

\subsection{Funding Stage and Layoff Severity}
Page 2 confirms that a company's funding stage heavily influences its layoff severity profile. Early-stage companies (e.g., Seed, Series A) exhibit higher \textit{layoff percentages} because they operate with limited cash reserves, fragile business models, and face existential threats when unable to raise new venture capital. Conversely, Late-stage or Post-IPO companies exhibit larger \textit{absolute layoff counts} because they employ tens of thousands of people and face relentless shareholder pressure to optimize profitability.

Therefore, funding stage is not merely a descriptive label; it is a proxy for company maturity, capital access, and structural resilience.

\subsection{Funds Raised and Restructuring Pressure}
The KPI for Total Funds Raised and the scatter plot contextualize layoffs against historical investment. High accumulated funding does not immunize a firm from layoffs; paradoxically, it often indicates that the firm previously scaled aggressively and later required severe right-sizing when growth expectations normalized.

With total funding reaching approximately \textbf{\$3.43 Trillion USD} across \textbf{4,411} companies, the scatter plot highlights critical outliers: several massively funded ``unicorn'' companies still executed large-scale layoffs. This proves that for many firms, layoffs were a strategic correction of over-hiring rather than a pure symptom of bankruptcy.

\subsection{Absolute vs. Relative Interpretation}
Ultimately, the dashboard's analytical strength lies in distinguishing between \textit{Total Layoffs} and \textit{Percentage Laid Off}. Total layoffs measure the sheer volume of affected workers, capturing the external macroeconomic shock to the labor market. Percentage laid off measures the internal severity of the restructuring. 

A multinational tech giant may record a staggering total layoff count while only reducing its workforce by 5\%. In contrast, a small startup might lay off just 50 people, but if that represents 80\% of its staff, it signals deep distress or imminent closure. This duality is central to interpreting the dashboard accurately.

\section{Research Questions and Visualization Mapping}
Table~\ref{tab:rq-mapping} explicitly maps the core research questions to the dataset variables, chosen visualization techniques, visual encodings, and their specific analytical purposes.

\begin{table}[H]
\centering
\caption{Mapping between research questions, dataset variables, visual encodings, and analytical purposes.}
\label{tab:rq-mapping}
\tiny
\renewcommand{\arraystretch}{1.3}
\setlength{\tabcolsep}{2pt}
\begin{tabular}{L{0.22\textwidth} L{0.13\textwidth} L{0.14\textwidth} L{0.24\textwidth} L{0.13\textwidth}}
\toprule
\textbf{Research Question} & \textbf{Dataset Variables} & \textbf{Visualization Technique} & \textbf{Visual Encoding} & \textbf{Analytical Purpose} \\
\midrule
How has the total number of layoffs changed over time? & \texttt{date}, \texttt{total\_laid\_off} & Area chart & x-axis = date; y-axis = \texttt{total\_laid\_off} & Identify temporal trends and macro peaks. \\
Which industries experienced the largest absolute layoffs? & \texttt{industry}, \texttt{total\_laid\_off} & Clustered bar chart & category = industry; bar length = \texttt{total\_laid\_off} & Compare sector-level impact. \\
Which countries are most affected by workforce reductions? & \texttt{country}, \texttt{total\_laid\_off}, Total Company & Choropleth map with tooltips & region = country; color intensity = \texttt{total\_laid\_off}; tooltip = Total Company & Analyze spatial concentration and scope. \\
How does funding stage relate to relative layoff severity? & \texttt{stage}, percentage laid off & Column chart & x-axis = stage; y-axis = mean percentage laid off & Evaluate distress across maturity stages. \\
Which specific companies contribute the most to the layoff volume? & \texttt{company}, \texttt{total\_laid\_off} & Clustered bar chart & category = company; bar length = \texttt{total\_laid\_off} & Identify dominant firm-level contributors. \\
Does high historical funding protect against large-scale layoffs? & \texttt{funds\_raised}, \texttt{total\_laid\_off} & Scatter plot & x = \texttt{funds\_raised}; y = \texttt{total\_laid\_off} & Examine the capital--layoff correlation and outliers. \\
\bottomrule
\end{tabular}
\normalsize
\end{table}

\section{Limitations}
While the dashboard provides a highly effective exploratory view of global technology layoffs, viewers must interpret the visualizations with an understanding of the following analytical and data-related limitations:

\begin{itemize}
    \item \textbf{Data Collection Bias:} The Kaggle Layoffs dataset relies heavily on media reports, press releases, and public disclosures. This inherently skews the data toward publicly traded technology giants and high-profile startups, while smaller private firms often conduct workforce reductions without public announcement.
    
    \item \textbf{The Disproportionate Weight of Megacorporations:} Absolute metrics, especially \texttt{total\_laid\_off}, are constrained by baseline company size. A few multinational tech giants can dominate aggregate layoff counts, industry bar charts, and the geographic choropleth map. This skew can mask severe internal distress among smaller startups if viewers consider only absolute volume. Therefore, total layoffs must be compared with percentage laid off to separate broad market impact from organizational severity.
    
    \item \textbf{Exploratory Scope and Non-Causal Associations:} The Power BI dashboard is designed for exploratory data analysis (EDA) and pattern recognition. It highlights correlations, such as which funding stages exhibit the highest relative layoff severity, but it cannot establish statistical causality. Variables like \texttt{stage} and \texttt{funds\_raised} are descriptive grouping dimensions, not proven triggers. The visualizations help generate hypotheses about macroeconomic pressure without proving causal mechanisms behind specific restructuring events.
\end{itemize}

\section{Conclusion}
This report has detailed the design, rationale, and implementation of an analytical Power BI dashboard tailored to investigate the global wave of layoffs. By rigorously applying foundational data visualization principles, the project successfully transformed a fragmented, tabular dataset into a cohesive analytical narrative. 

The dual-page architecture systematically guides the viewer from macro to micro perspectives. While Page 1 establishes the global context, identifying the precise temporal peaks of the layoff wave, mapping its geographic epicenter, and highlighting the disproportionate vulnerability of specific sectors like Retail. Page 2 deepens this analysis by probing the structural characteristics of the firms themselves, revealing how a company's funding maturity and accumulated capital dictate its restructuring behavior.

Ultimately, this project demonstrates that when raw data is strictly prepared, modeled, and visually encoded with precision, it transitions from a simple historical record into a powerful diagnostic instrument. It empowers analysts, researchers, and strategic decision-makers to navigate complex economic signals, evaluate firm-level vulnerabilities, and extract actionable insights from global labor-market fluctuations.
\end{document}
